\begin{recipe}{Parmezaan-Knoflook-Rozemarijn Kipfilet in Bakpapier}
\kicker{Hoofdgerecht · vlees · doordeweeks}
\index[register]{Parmezaan-Knoflook Kipfilet}
\index[register]{Kipfilet}
\index[register]{Parmezaanse kaas}
\index[register]{Knoflook}

\meta{10-15 minuten \dvd Voor 2 personen \dvd Snel}

\lettrine{J}{aren} geleden at ik wel eens Maggi Mals — die verpakkings-kipfilets met kaas en kruiden — en vond het concept briljant: malse, kruimige kipfilet klaar in minuten. Dit is mijn versie om het zelf te maken: veel simpeler, veel goedkoper, en minstens zo sappig.

\blockrule{Ingrediënten}
\begin{ingredients}
  \ing{2}{kipfilets (ca. 180 g per stuk)}
  \ing{5 el}{geraspte Parmezaanse kaas}
  \ing{3 teentjes}{knoflook, fijngehakt (of 1 tl knoflookpoeder)}
  \ing{1½ tl}{gedroogde rozemarijn (of 1 el verse, fijngehakt)}
  \ing{1½ tl}{olijfolie}
  \ingb{zout naar smaak}
  \ingb{versgemalen zwarte peper}
  \ingb{bakpapier}
\end{ingredients}

\blockrule{Bereiding}
\tip{Middelhoog vuur is ideaal voor een mooi kaaslaagje zonder aanbranden.}
\begin{steps}
  \item Leg een kipfilet op een snijplank en snij hem horizontaal bijna helemaal doormidden, zodat je hem open kunt klappen als een boek. Herhaal met de tweede filet.
  \item Leg het gevlinderde stuk tussen 2 vellen bakpapier en sla het voorzichtig plat met een deegroller, vleeshamer of de bodem van een zware pan tot een gelijkmatige dikte van ca. 1 cm.
  \item Haal het bovenste vel bakpapier eraf. Bestrooi de kip royaal met Parmezaanse kaas, zout, peper, knoflook en rozemarijn. Druk alles goed aan en druppel ½-1 tl olijfolie erover.
  \item Leg het tweede vel bakpapier erop en druk stevig aan, zodat de kip helemaal bedekt is.
  \item Verhit een koekenpan op middelhoog vuur. Laat de pan 1 minuut goed warm worden en leg het bakpapier-pakketje erin.
  \item Bak 5-7 minuten per kant, keer voorzichtig om met een spatel. De kip is gaar bij een kerntemperatuur van 75°C (of als het vlees wit is en het sap helder loopt).
  \item Haal uit de pan, laat 2-3 minuten rusten en verwijder het bakpapier. Snij de kip in plakken of serveer heel. Lekker met salade, groenten, pasta, rijst of aardappelen.
\end{steps}
\end{recipe}
